Capitulo: Traps e Chamadas de Sistema
Ha tres tipos de eventos que fazem a CPU interromper a execucao normal das instrucoes e transferir o
controle para um codigo especial que trata o evento. Um deles e uma chamada de sistema, quando um
programa em modo usuario executa a instrucao `ecall` para solicitar algo ao kernel. Outro e uma excecao:
quando uma instrucao faz algo ilegal, como dividir por zero ou acessar um endereco invalido. O terceiro tipo
e uma interrupcao de dispositivo, como quando o hardware do disco conclui uma leitura ou escrita.
Este livro usa o termo "trap" como uma referencia geral a esses eventos. A ideia e que o codigo que estava
sendo executado no momento da trap possa continuar depois, sem saber que foi interrompido. Traps devem
ser transparentes, especialmente as interrupcoes de dispositivos. A sequencia usual e: a trap transfere o
controle ao kernel; o kernel salva os registradores e o estado; executa o codigo apropriado (por exemplo,
chamada de sistema ou driver); restaura o estado salvo e retorna de onde parou.
O xv6 trata todas as traps no kernel. Isso e natural para chamadas de sistema e necessario para
interrupcoes, ja que o kernel e responsavel por controlar os dispositivos e garantir o isolamento entre
processos. Tambem e apropriado para excecoes: o xv6 termina qualquer processo que cause uma excecao.
O tratamento de traps no xv6 ocorre em quatro fases: acoes de hardware da CPU RISC-V, instrucoes em
assembly que preparam o caminho para o codigo C do kernel, uma funcao C que decide como tratar a trap,
e a rotina correspondente (chamada de sistema ou driver).
Cada CPU RISC-V tem registradores de controle (como `stvec`, `sepc`, `scause`, `sscratch`, `sstatus`) que
participam do processo de tratamento de traps. Quando ocorre uma trap, o hardware realiza uma sequencia
de etapas: desativa interrupcoes, salva o contador de programa, registra a causa, muda para o modo
supervisor e transfere o controle para o endereco em `stvec`.
No xv6, traps vindas do espaco do usuario seguem o fluxo: `uservec` `usertrap` `usertrapret` `userret`.
Uma limitacao do RISC-V e que ele nao troca a tabela de paginas automaticamente ao entrar no modo
supervisor. Para contornar isso, o xv6 usa uma pagina especial chamada *trampoline*, mapeada tanto na
tabela de paginas do usuario quanto do kernel.
O tratador `uservec` salva os 32 registradores do usuario no `trapframe` e muda a tabela de paginas para a
do kernel. Em seguida, `usertrap` identifica a causa da trap: se for chamada de sistema, chama `syscall`; se
for interrupcao, chama `devintr`; se for excecao, mata o processo.
Na volta, `usertrapret` prepara os registradores e chama `userret`, que restaura os registradores e executa
`sret` para voltar ao usuario.
Chamadas de sistema como `exec()` sao iniciadas com `ecall`. O numero da chamada vai em `a7` e os
argumentos nos registradores. O kernel busca esses valores no `trapframe` e executa a funcao
correspondente, salvando o valor de retorno em `a0`.
Para acessar ponteiros passados pelo usuario, o kernel usa funcoes como `copyin`, `copyout` e `copyinstr`,
que validam os enderecos e copiam dados entre os espacos de memoria do kernel e do usuario.
Traps originadas no kernel seguem um caminho diferente. O handler `kernelvec` salva os registradores na
pilha do kernel e chama `kerneltrap`, que trata interrupcoes e, se necessario, executa `yield`. Depois, os
registradores sao restaurados e a execucao e retomada com `sret`.
Em sistemas reais, traps e excecoes como falhas de pagina sao usadas para implementar recursos como
*copy-on-write fork*, *lazy allocation* (alocacao preguicosa), *demand paging* (paginacao sob demanda), e
*paging to disk* (paginacao em disco). Esses mecanismos aumentam a eficiencia e reduzem o uso de
memoria, permitindo que sistemas operacionais modernos escalem melhor.
Alem disso, recursos como pilhas autoexpansiveis e arquivos mapeados em memoria (`mmap`) sao
suportados atraves de excecoes de pagina. Sistemas de producao tambem fornecem chamadas como
`mmap`, `munmap`, `sigaction`, `mlock`, `madvise`, para controle avancado da memoria pelo usuario.
Por fim, embora estruturas como `trampoline` e `trapframe` possam parecer complexas, elas sao uma
consequencia do design minimalista da RISC-V, que deixa para o software grande parte da logica de
tratamento de traps, permitindo flexibilidade e desempenho.
Exercicios sugeridos incluem implementar alocacao preguicosa, fork com copia-sob-escrita, `mmap`, e
simplificar o uso de `TRAMPOLINE` e `TRAPFRAME`.
